% !PRINTED 357
\weberpage{357}
This equation shows, however, that the curve of degree $14$:

\weq{19}{\chi(x) = 0,}

which is entirely independent of the $a$, passes through the points of contact of the double tangents, but through no other point of the curve $f(x)$.

The function $\chi(x)$ depends rationally on the coefficients of the equation $f(x)$. There are infinitely many different such curves, among which covariants of $f(x)$ are also found. As \gesperrt{Hesse} has shown, they can be expressed in a simple way by determinants.

Here we draw from this the conclusion:

2.\ \gesperrt{A curve of the fourth order without a singular point has 28 double tangents.}

But one objection to this conclusion must still be met. It would be conceivable that the curve $\chi$ touches the curve $f$, or that the curve $f$ passes through singular points of the curve $\chi$. Then the number of points of intersection, and possibly also the number of double tangents, would be diminished, so that our mode of inference really teaches only that a curve of the fourth order cannot have more than 28 double tangents. But this objection will dispose of itself through the following considerations, in that we shall become acquainted with forms of the equation of the curve in which the existence of 28 distinct double tangents is evident.

\begin{center}
\S.\,96.

\gesperrt{The Steiner Complexes.}
\end{center}

If $x_1 = 0$ is the equation of any double tangent of the curve of the fourth order $f = 0$, which we call for short the double tangent $x_1$, then, when one sets $x_1 = 0$, $f$ must pass into a square. From this it results that $f$ must be of the form:

\weq{1}{f = x_1 V - u^2,}

in which $V$ is a cubic, $u$ a quadratic form. But the function $f$ can be given the shape (1) in a triply infinite number of ways. For if $p$ denotes an arbitrary
%CONT>>
